\chapter{Release 3 --- Pipeline CI/CD}

\section*{Introduction}
\addcontentsline{toc}{section}{Introduction}
Cette release vise à automatiser l'intégration et la livraison continues du projet. Elle couvre la mise en place du pipeline Jenkins, l'analyse de qualité avec SonarQube, le scan de sécurité avec Trivy, le packaging avec Helm et le déploiement continu avec ArgoCD.

%======================================================================
% SPRINT 6
%======================================================================
\section{Sprint 6 --- Pipeline CI/CD et déploiement continu}

\subsection{Objectif du sprint}
Mettre en place un pipeline d'intégration et de livraison continues complet, depuis le commit du code jusqu'au déploiement automatique sur le cluster EKS.

\subsection{Sprint Backlog}

\begin{table}[H]
\centering
\renewcommand{\arraystretch}{1.2}
\begin{tabular}{|c|p{9cm}|c|c|}
\hline
\textbf{ID} & \textbf{Tâche} & \textbf{Priorité} & \textbf{Points} \\
\hline
6.1 & Configuration du serveur Jenkins (Dockerfile custom) & Must & 5 \\
\hline
6.2 & Pipeline multi-étapes : build Maven, tests unitaires & Must & 8 \\
\hline
6.3 & Intégration SonarQube pour l'analyse de qualité du code & Should & 5 \\
\hline
6.4 & Scan de sécurité des images avec Trivy & Should & 3 \\
\hline
6.5 & Push automatisé des images vers AWS ECR & Must & 5 \\
\hline
6.6 & Configuration d'ArgoCD et Helm pour le déploiement GitOps & Must & 8 \\
\hline
\end{tabular}
\caption{Sprint Backlog --- Sprint 6}
\label{tab:sprint6_backlog}
\end{table}

\subsection{Architecture du pipeline CI/CD}

Le pipeline CI/CD suit un modèle GitOps. Le processus complet se déroule en deux phases :

\begin{enumerate}
    \item \textbf{Intégration continue (CI)} --- gérée par Jenkins :
    \begin{itemize}
        \item Checkout du code source depuis GitHub.
        \item Compilation avec Maven et exécution des tests unitaires.
        \item Analyse de la qualité du code avec SonarQube.
        \item Construction des images Docker.
        \item Scan de sécurité des images avec Trivy.
        \item Push des images vers AWS ECR.
        \item Mise à jour des tags d'images dans le dépôt de configuration GitOps.
    \end{itemize}
    \item \textbf{Livraison continue (CD)} --- gérée par ArgoCD :
    \begin{itemize}
        \item ArgoCD surveille le dépôt \texttt{billing-app-config} (repo GitOps).
        \item Détection automatique des changements de tags d'images.
        \item Synchronisation automatique des manifests Kubernetes sur le cluster EKS.
    \end{itemize}
\end{enumerate}

\begin{figure}[H]
    \centering
    \fbox{\parbox{0.85\textwidth}{\centering\vspace{3cm}\textbf{[Diagramme du pipeline CI/CD]}\\\textit{GitHub $\rightarrow$ Jenkins (Build, Test, Scan, Push) $\rightarrow$ ECR $\rightarrow$ Config Repo $\rightarrow$ ArgoCD $\rightarrow$ EKS}\vspace{3cm}}}
    \caption{Architecture du pipeline CI/CD}
    \label{fig:cicd_pipeline}
\end{figure}

\subsection{Serveur Jenkins}
Le serveur Jenkins est déployé sur une instance EC2 dédiée. Son image Docker est personnalisée pour inclure tous les outils nécessaires :

\begin{lstlisting}[language=bash, caption={Dockerfile Jenkins personnalisé}, label={lst:jenkins_dockerfile}]
FROM jenkins/jenkins:lts
USER root
RUN apt-get update && \
    apt-get install -y docker.io git curl maven && \
    apt-get clean
# Trivy
RUN curl -sfL https://raw.githubusercontent.com/aquasecurity/
    trivy/main/contrib/install.sh | sh -s -- -b /usr/local/bin
# Helm
RUN curl https://raw.githubusercontent.com/helm/helm/
    main/scripts/get-helm-3 | bash
# AWS CLI
RUN curl "https://awscli.amazonaws.com/awscli-exe-
    linux-x86_64.zip" -o "awscliv2.zip" && \
    unzip awscliv2.zip && ./aws/install
\end{lstlisting}

\subsection{Pipeline Jenkinsfile}
Le pipeline est défini dans un fichier \texttt{Jenkinsfile} déclaratif à la racine du projet.

\begin{figure}[H]
    \centering
    \fbox{\parbox{0.85\textwidth}{\centering\vspace{2.5cm}\textbf{[Capture d'écran --- Pipeline Jenkins]}\\\textit{Vue des stages : Checkout $\rightarrow$ Build $\rightarrow$ SonarQube $\rightarrow$ Docker $\rightarrow$ Trivy $\rightarrow$ ECR $\rightarrow$ Update Config}\vspace{2.5cm}}}
    \caption{Pipeline Jenkins avec les différentes étapes}
    \label{fig:jenkins_pipeline}
\end{figure}

\subsection{Analyse de qualité --- SonarQube}
SonarQube est utilisé pour l'analyse statique du code. Il détecte les bugs, les vulnérabilités et les mauvaises pratiques de code (code smells).

\begin{figure}[H]
    \centering
    \fbox{\parbox{0.85\textwidth}{\centering\vspace{2.5cm}\textbf{[Capture d'écran --- Dashboard SonarQube]}\vspace{2.5cm}}}
    \caption{Résultats d'analyse SonarQube}
    \label{fig:sonarqube}
\end{figure}

\subsection{Scan de sécurité --- Trivy}
Trivy est un scanner de vulnérabilités pour les images Docker. Il est exécuté dans le pipeline Jenkins après la construction de chaque image.

\subsection{Déploiement GitOps avec ArgoCD}
ArgoCD surveille le dépôt de configuration \texttt{billing-app-config} et synchronise automatiquement les manifests Kubernetes sur le cluster EKS.

\begin{figure}[H]
    \centering
    \fbox{\parbox{0.85\textwidth}{\centering\vspace{2.5cm}\textbf{[Capture d'écran --- ArgoCD Dashboard]}\\\textit{Applications synchronisées avec leur statut de santé}\vspace{2.5cm}}}
    \caption{Dashboard ArgoCD}
    \label{fig:argocd}
\end{figure}

\subsection{Packaging avec Helm}
Helm est utilisé pour le packaging des déploiements Kubernetes en production. Les Helm charts permettent de paramétrer les déploiements (image tags, replicas, variables d'environnement) via un fichier \texttt{values.yaml}.

\section*{Conclusion}
\addcontentsline{toc}{section}{Conclusion}
Au terme de cette release, le pipeline CI/CD est pleinement opérationnel. Chaque commit déclenche automatiquement la compilation, les tests, l'analyse de qualité, le scan de sécurité et le déploiement sur le cluster EKS via ArgoCD. Le chapitre suivant abordera la mise en place du monitoring.

\newpage
